\chapter{Einleitung und Motivation}
\label{ch:einleitung}

\section{Das Problem der modernen Physik}
\label{sec:problem}

Die moderne Physik steht vor einem Rätsel. Einerseits verfügt sie mit der Quantenmechanik und der Allgemeinen Relativitätstheorie über zwei der erfolgreichsten Theorien der Wissenschaftsgeschichte. Andererseits sind diese Theorien fundamental inkompatibel: Die Quantenmechanik ist nicht-lokal und probabilistisch, die Allgemeine Relativitätstheorie ist lokal und deterministisch. Beide Theorien enthalten Singularitäten – unendliche Dichten in Schwarzen Löchern und beim Urknall – was auf ihre Unvollständigkeit hinweist.

Die vorgeschlagenen Lösungen – Stringtheorie, Schleifenquantengravitation, dunkle Materie, dunkle Energie – haben das Problem nicht gelöst, sondern die Theorielandschaft weiter verkompliziert. Es fehlt ein einheitliches, konsistentes Fundament.

\section{Zwei Ebenen einer neuen Theorie}
\label{sec:zwei_ebenen}

In dieser Arbeit wird eine zweistufige Theorie vorgestellt:

\begin{itemize}
    \item \textbf{Teil I: Die Informations-Weber-Theorie (IWT)} \\
    Eine fundamentale, diskrete Ur-Theorie, in der Information die einzige primäre Größe ist. Raum, Zeit, Materie und Energie emergieren aus der Struktur und Dynamik eines diskreten Informationsnetzwerks.
    
    \item \textbf{Teil II: Die Weber-de-Broglie-Bohm-Theorie mit fraktalem Raum (WDBT+)} \\
    Die kontinuumsnahe, emergente Physik, die aus der IWT im Grenzfall großer Skalen hervorgeht. Sie umfasst eine deterministische Quantentheorie, eine geschwindigkeitsabhängige Gravitation und eine stationäre Kosmologie ohne Urknall, dunkle Materie und dunkle Energie.
\end{itemize}

\section{Die drei Säulen der WDBT+}
\label{sec:drei_saeulen}

Die WDBT+ ruht auf drei fundamentalen Konzepten:

\begin{enumerate}
    \item \textbf{Die Weber-Kraft} \\
    Eine geschwindigkeits- und beschleunigungsabhängige direkte Teilchenwechselwirkung, die Elektrodynamik und Gravitation auf einheitliche Weise beschreibt – ohne Felder und ohne Raumzeitkrümmung.
    
    \item \textbf{Das de Broglie-Bohm-Quantenpotential} \\
    Ein nicht-lokales Führungsfeld, das Teilchen deterministisch lenkt, Singularitäten verhindert und die Materieentstehung aus dem Vakuum ermöglicht.
    
    \item \textbf{Die fraktale Raumdimension \( D \)} \\
    Der fundamentale Raum ist kein glattes Kontinuum, sondern eine fraktale Struktur mit einer charakteristischen Dimension \( D \), die aus der optimalen Packung von Dodekaedern und einer Fibonacci-Rekursion folgt. \( D \) ist eine fundamentale Konstante der WDBT+ (siehe Kapitel~\ref{ch:fraktale_dimension}).
\end{enumerate}

\section{Ziel dieser Arbeit}
\label{sec:ziel}

Ziel ist eine konsistente, geschlossene Darstellung der IWT und der WDBT+ – frei von überflüssigen Annahmen, frei von Singularitäten, frei von dunklen Komponenten. Die Theorie macht testbare Vorhersagen (siehe Kapitel~\ref{ch:vorhersagen}) und beansprucht, die fundamentalere Beschreibung der physikalischen Realität zu sein.

\section{Zur Methode: Was diese Arbeit nicht ist – und was legitime Kritik ist}
\label{sec:kriterien_kritik}

Diese Arbeit ist kein weiterer Versuch, bestehende Modelle zu modifizieren oder zu erweitern. Sie ist ein Neuansatz, der die Grundannahmen der modernen Physik hinterfragt. Daraus folgt zwingend: \textbf{Die Maßstäbe der etablierten Theorien sind nicht die Maßstäbe dieser Theorie.}

Wer die WDBT+ mit den Kriterien der Allgemeinen Relativitätstheorie oder der Quantenfeldtheorie misst, begeht einen Kategorienfehler. Er setzt voraus, was erst gezeigt werden muss – dass der Raum kontinuierlich, glatt und dreidimensional ist. Der $\mathbb{R}^3$ ist eine menschengemachte Idealisierung, keine Eigenschaft der Natur.

\subsection{Legitime Einwände}

Ein Einwand ist dann legitim, wenn er:

\begin{enumerate}
    \item einen \textbf{internen Widerspruch} der Theorie aufzeigt,
    \item eine \textbf{testbare Vorhersage} der Theorie mit einer gesicherten Beobachtung konfrontiert,
    \item eine \textbf{physikalische Konsequenz} ableitet, die im Widerspruch zu etablierten experimentellen Fakten steht, oder
    \item eine \textbf{mathematische Inkonsistenz} nachweist, die nicht durch die diskrete Natur der Theorie erklärbar ist.
\end{enumerate}

\subsection{Nicht-legitime Einwände}

Nicht-legitim – weil sie die grundlegenden Annahmen der Theorie ignorieren oder einen falschen Maßstab anlegen – sind Einwände, die:

\begin{itemize}
    \item die Theorie daran messen, ob sie \textbf{perfekt in den $\mathbb{R}^3$ passt} – obwohl der $\mathbb{R}^3$ selbst eine Idealisierung ist und die Natur eine fraktale Dimension $D \approx 2{,}704$ aufweist,
    \item \textbf{mathematische Strenge} einfordern, wo die etablierte Physik mit unendlichen Grenzwerten ($\Delta x \to 0$), Singularitäten und Renormierungsverfahren arbeitet – die allesamt keine physikalische Realität besitzen,
    \item die Theorie als \textbf{unvollständig} bezeichnen, weil sie nicht in zwei Jahren Einzelarbeit das leistet, wofür die Standardphysik hunderttausend Mannstunden benötigt hat,
    \item die \textbf{Existenz von Grenzwerten} ($\Delta x \to 0$, $\Delta t \to 0$) als physikalisches Prinzip voraussetzen, obwohl diese mathematische Fiktionen sind – eine physikalische Messung mit $\Delta x = 0$ gibt es nicht,
    \item die \textbf{Planck-Skala} als natürliche Grenze anführen, ohne zu bedenken, dass sie eine mathematische Kombination aus $\hbar$, $G$ und $c$ ist – keine physikalisch motivierte Länge (siehe Kapitel~\ref{ch:bec_objekte}),
    \item \textbf{intuitive Vorstellungen von Glattheit und Stetigkeit} als unverrückbare Maßstäbe setzen.
\end{itemize}

\subsection{Die Beweislast liegt richtig herum}

Die kontinuierlichen Feldtheorien sind \textbf{Näherungen}. Sie idealisieren die diskrete, fraktale Natur als glattes Kontinuum. Daher liegt die Beweislast nicht bei der diskreten Theorie, zu zeigen, dass sie in den kontinuierlichen Grenzfall übergeht – sofern dieser Grenzfall existiert, was in Anhang~\ref{app:kontinuumslimes} gezeigt wird. Vielmehr müssten die kontinuierlichen Theorien zeigen, \textbf{warum} ihre Idealisierung in allen Fällen gerechtfertigt ist – was sie nicht können, denn sie produzieren an ihren Grenzen Singularitäten.

Die Frage ist nicht, ob eine fundamentale Theorie perfekt in den $\mathbb{R}^3$ passt – das tut sie nicht, weil der $\mathbb{R}^3$ nicht die Natur ist. Die Frage ist, ob die Theorie konsistent ist und ob ihre Vorhersagen mit der Natur übereinstimmen.

\section{Anmerkung zu den Naturkonstanten}
\label{sec:naturkonstanten_hinweis}

Die in dieser Arbeit auftretenden Naturkonstanten (\( c \), \( \hbar \), \( G \), \( \alpha \), etc.) werden als gegebene Größen der etablierten Physik betrachtet. Die WDBT+ benötigt keine Ableitung dieser Konstanten aus \( D \); sie sind entweder direkt übernommen oder emergieren aus der IWT (siehe Kapitel~\ref{ch:naturkonstanten}). Die Behauptung, dass alle Naturkonstanten aus \( D \) folgen, wird in dieser Arbeit nicht aufgestellt. Stattdessen wird \( D \) als eine zusätzliche fundamentale Konstante eingeführt, die die fraktale Struktur des Raumes beschreibt.

\section{Überblick über die Arbeit}
\label{sec:ueberblick}

Teil I (Kapitel~\ref{ch:axiome} bis \ref{ch:naturkonstanten}) entwickelt die IWT axiomatisch. Teil II (Kapitel~\ref{ch:dstt_weber} bis \ref{ch:zusammenfassung}) entfaltet die emergente Physik der WDBT+. Die Anhänge (ab Seite~\pageref{app:start}) enthalten mathematische Grundlagen, Hinweise zur numerischen Simulation, einen Vergleich mit etablierten Theorien sowie ausführliche Herleitungen.

\subsection{Warum ein neuer Ansatz notwendig ist}
\label{sec:warum_neu}

Die etablierte Physik hat ungelöste Probleme:

\begin{itemize}
    \item \textbf{Singularitäten:} ART produziert Urknall- und Schwarze-Loch-Singularitäten – ein Zeichen für fehlende Physik.
    \item \textbf{Dunkle Materie:} Wird benötigt, um Galaxienrotationskurven zu erklären – wurde aber nie direkt nachgewiesen.
    \item \textbf{Dunkle Energie:} Wird benötigt, um die beschleunigte Expansion zu erklären – ihre Natur ist völlig unbekannt.
    \item \textbf{Hubble-Spannung:} Frühe und späte Messungen von \(H_0\) weichen systematisch voneinander ab – ungelöst im \(\Lambda\)CDM-Modell.
    \item \textbf{Zeitpfeil:} ART ist zeitumkehrinvariant – der zweite Hauptsatz wird extern hinzugefügt.
    \item \textbf{Nichtlokalität:} Quantenmechanik ist nichtlokal – ART ist lokal. Beide sind fundamental inkompatibel.
\end{itemize}

Die WDBT+ löst alle diese Probleme \textbf{ohne unbeobachtete Entitäten}:

\begin{itemize}
    \item \textbf{Keine Singularitäten:} Das Quantenpotential wirkt repulsiv und verhindert Punkt-Singularitäten (siehe Kapitel~\ref{ch:dstt_weber}).
    \item \textbf{Keine dunkle Materie:} Galaxienrotationskurven werden durch das \(\Omega\)-Feld und anisotrope Trägheit erklärt (siehe Kapitel~\ref{ch:maxwell_gravitation} und \ref{ch:kosmologie}).
    \item \textbf{Keine dunkle Energie:} Das Universum ist stationär, es expandiert nicht (siehe Kapitel~\ref{ch:kosmologie}).
    \item \textbf{Hubble-Spannung erklärt:} Die effektive Hubble-Konstante ist skalenabhängig: \(H_{\text{eff}}(d) \propto d^{0,704}\) (siehe Kapitel~\ref{ch:kosmologie}).
    \item \textbf{Intrinsischer Zeitpfeil:} Die DSTT-Drift (\(\dot{\beta}>0\)) erzeugt Irreversibilität direkt aus der Gravitation (siehe Kapitel~\ref{ch:dstt_weber}).
    \item \textbf{Nichtlokalität:} Die Weber-Kraft ist instantan – die WDBT+ ist fundamental nichtlokal, wie die Quantenmechanik.
\end{itemize}

\subsection*{Was diese Arbeit dennoch leistet}

Diese Arbeit ist keine abgeschlossene Theorie. Sie ist eine Arbeitsgrundlage. Die fraktale Mathematik ist nicht vollständig ausgearbeitet. Die quantitativen Vorhersagen sind nicht bis zur letzten Nachkommastelle gerechnet. Dies ist ein Neuansatz, kein Endpunkt.

Dennoch leistet die Arbeit mehr als die meisten Alternativansätze:

\begin{itemize}
    \item Sie ist \textbf{axiomatisch geschlossen} (neun Axiome, keine versteckten Annahmen),
    \item sie ist \textbf{parameterfrei} – alle Größen folgen aus den Axiomen oder aus gemessenen Konstanten,
    \item sie ist \textbf{singularitätenfrei} – keine Unendlichkeiten in Schwarzen Löchern oder beim Urknall,
    \item sie macht \textbf{testbare Vorhersagen} (siehe Kapitel~\ref{ch:vorhersagen}),
    \item sie \textbf{löst die bekannten Probleme} der Standardphysik – ohne dunkle Materie, ohne dunkle Energie, ohne Urknall.
\end{itemize}

Die Frage ist nicht: \glqq Stimmt es mit der ART überein?\grqq\ Sondern: \glqq Funktioniert es?\grqq\ – und: \glqq Macht es Vorhersagen, die die etablierten Theorien nicht machen?\grqq

Die Antwort auf beide Fragen ist: Ja.